% !TeX root = ../sections/02_exercises.tex

\subsection{3-Review-1}
\begin{exercise}
	\(R=\mathbb{R}[x]\) とし，\(V=\mathbb{R}^2\) とする。
	\(2\) 次行列
	\[
	A=
	\begin{pmatrix}
		3 & 1\\
		0 & 2
	\end{pmatrix}
	\]
	を用いて，\(\mathbb{R}[x]\) の \(V=\mathbb{R}^2\) への作用を
	\[
	x
	\begin{pmatrix}
		a\\
		b
	\end{pmatrix}
	=
	A
	\begin{pmatrix}
		a\\
		b
	\end{pmatrix}
	\]
	で定義する。
	
	\begin{enumerate}
		\item 次を計算せよ。
		\[
		\text{\rm (i)}\ 
		(2x+3)
		\begin{pmatrix}
			4\\
			5
		\end{pmatrix}
		\qquad
		\text{\rm (ii)}\ 
		(x^2-2)
		\begin{pmatrix}
			3\\
			-2
		\end{pmatrix}
		\]
		
		\item
		\[
		W=
		\left\{
		\begin{pmatrix}
			a\\
			0
		\end{pmatrix}
		\ \middle|\ 
		a\in\mathbb{R}
		\right\}
		\]
		は \(V\) の \(\mathbb{R}[x]\)-部分加群であることを示せ。
		
		\item
		\(W\) は既約 \(\mathbb{R}[x]\)-加群であることを示せ。
		
		\item
		剰余加群
		\[
		V/W=
		\left\{
		\left[
		\begin{pmatrix}
			0\\
			b
		\end{pmatrix}
		\right]
		\ \middle|\ 
		b\in\mathbb{R}
		\right\}
		\]
		において，次の計算をせよ。
		\[
		\text{\rm (i)}\ 
		(3x-2)
		\left[
		\begin{pmatrix}
			0\\
			3
		\end{pmatrix}
		\right]
		\qquad
		\text{\rm (ii)}\ 
		(2x^2+3)
		\left[
		\begin{pmatrix}
			0\\
			5
		\end{pmatrix}
		\right]
		\]
	\end{enumerate}
\end{exercise}

\begin{background}
	この問題では，ベクトル空間 \(V=\mathbb{R}^2\) を
	\(\mathbb{R}[x]\)-加群とみなす。
	ここで \(x\) の作用は行列 \(A\) によって定められているので，
	多項式 \(f(x)\in\mathbb{R}[x]\) の作用は
	\[
	f(x)v=f(A)v
	\]
	として計算される。
	
	特に，
	\[
	(2x+3)v=(2A+3I)v,
	\qquad
	(x^2-2)v=(A^2-2I)v
	\]
	である。
	
	また，\(W\subset V\) が \(\mathbb{R}[x]\)-部分加群であるとは，
	\(W\) が加法について部分群であり，さらに任意の
	\(f(x)\in\mathbb{R}[x]\) と \(w\in W\) に対して
	\[
	f(x)w\in W
	\]
	が成り立つことをいう。
	
	剰余加群 \(V/W\) では，二つのベクトルはその差が \(W\) に属するとき
	同じ元とみなされる。
	したがって，第1成分の差は \(W\) に吸収され，第2成分が本質的に残る。
\end{background}

\begin{strategy}
	まず，\(\mathbb{R}[x]\)-加群の作用を行列の計算に翻訳する。
	つまり，\(x\) は \(A\) として作用するので，
	多項式 \(f(x)\) は行列 \(f(A)\) として作用する。
	
	\begin{enumerate}
		\item[(1)]
		\((2x+3)v\) は \((2A+3I)v\) として計算する。
		同様に，\((x^2-2)v\) は \((A^2-2I)v\) として計算する。
		
		\item[(2)]
		\(W\) が通常の加法について部分群であることを確認した後，
		\(xW\subset W\) を示す。
		\(\mathbb{R}[x]\) は \(x\) と定数から作られるので，
		\(x\) の作用で \(W\) が保たれれば，
		任意の多項式の作用でも \(W\) が保たれる。
		
		\item[(3)]
		\(W\) の非零 \(\mathbb{R}[x]\)-部分加群 \(U\) を任意に取る。
		\(U\) に非零ベクトル
		\[
		\begin{pmatrix}
			a\\
			0
		\end{pmatrix}
		\qquad (a\neq 0)
		\]
		が入っていれば，定数倍の作用により
		\[
		\begin{pmatrix}
			1\\
			0
		\end{pmatrix}
		\]
		も入る。
		したがって \(U=W\) となることを示せばよい。
		
		\item[(4)]
		剰余加群では第1成分は \(W\) によって消える。
		そのため，
		\[
		A
		\begin{pmatrix}
			0\\
			b
		\end{pmatrix}
		=
		\begin{pmatrix}
			b\\
			2b
		\end{pmatrix}
		\]
		を計算した後，第1成分を無視して，
		同値類として整理する。
	\end{enumerate}
\end{strategy}